\documentclass[10pt,aspectratio=169]{beamer}

% --- pdflatex-safe font & micro-typography ---
\usepackage[T1]{fontenc}
\usepackage[utf8]{inputenc}
\usepackage{lmodern}
\usepackage{microtype}

% --- math & symbols ---
\usepackage{amsmath,amssymb,mathtools,bm}

% --- hyperlinks (Beamer already loads hyperref; this is safe) ---
\hypersetup{hidelinks}

% --- theme ---
\usetheme[numbering=fraction]{metropolis}

\usepackage{pgfpages} % for notes layouts

% --- bibliography ---
\usepackage[round,authoryear]{natbib}

% --- Choose ONE of these toggles when compiling ---
% \setbeameroption{hide notes}                        % audience version (default)
% \setbeameroption{show notes}                        % notes printed under each slide
% \setbeameroption{show only notes}                   % notes-only PDF (for printing)
% \setbeameroption{show notes on second screen=right} % slide + notes side-by-side

% --- shortcuts ---
\newcommand{\E}{\mathbb{E}}
\newcommand{\R}{\mathbb{R}}
\newcommand{\btheta}{\boldsymbol{\theta}}
\newcommand{\bx}{\boldsymbol{x}}

% --- title info ---
\title{Progress meeting}
\subtitle{Concentration inequalities}
\author{Shreyas Kalvankar}
\institute{TU Delft}
\date{\today}

\begin{document}

\maketitle
% -----------------------
\begin{frame}{Recap}
	\small
	We have three concentration lemmas for the first \(K\) Fourier modes:
	\begin{itemize}
		\item \textbf{Lemma A:} sampled Fourier features are almost orthogonal,
		\item \textbf{Lemma B:} the compressed operator is close to its expectation,
		\item \textbf{Lemma C:} the empirical operator acts approximately like the population eigen-operator on these modes.
	\end{itemize}

	\begin{block}{Question}
		If we fix a target accuracy \(\varepsilon\), confidence \(1-\delta\), and mode cutoff \(K\), what sample size \(n\) do the bounds actually require?
	\end{block}

	\note{This section is not about sharp constants; it is about honestly checking what the current concentration proofs imply numerically.}
\end{frame}

% -----------------------
\begin{frame}{Numerical choices for the back-of-the-envelope calculation}
	\small
	We plug in one concrete set of values:
	\[
		\varepsilon = 0.05,\qquad \delta = 0.05,\qquad K=4,\qquad d=2K+1=9.
	\]

	For the kernel bound we use
	\[
		\kappa=\|\Theta\|_\infty=\Theta(0)=\frac{3}{2},
	\]
	which is the worst-case bound coming from the closed-form NTK on \(S^1\).

	\begin{block}{Interpretation}
		\begin{itemize}
			\item \(\varepsilon=0.05\): target entrywise error level,
			\item \(\delta=0.05\): \(95\%\) confidence,
			\item \(K=4\): first 4 Fourier modes only.
		\end{itemize}
	\end{block}

	\note{We are deliberately using the exact constants from the lemmas, with no optimization.}
\end{frame}

% -----------------------
\begin{frame}{Lemma A: Gram concentration}
	\small
	The bound was
	\[
		\mathbb P\!\left(\left\|\frac1n\Phi^\top\Phi-I_d\right\|_{\max}\ge \varepsilon\right)
		\le 2d^2\exp\!\left(-\frac{n\varepsilon^2}{8}\right).
	\]

	To make the failure probability at most \(\delta\), it is enough that
	\[
		2d^2\exp\!\left(-\frac{n\varepsilon^2}{8}\right)\le \delta,
	\]
	which gives
	\[
		\boxed{
			n \ge \frac{8}{\varepsilon^2}\log\!\left(\frac{2d^2}{\delta}\right).
		}
	\]

	Plugging in \(\varepsilon=0.05\), \(d=9\), \(\delta=0.05\):
	\[
		n \ge \frac{8}{0.05^2}\log\!\left(\frac{2\cdot 9^2}{0.05}\right)
		\approx 25{,}867.
	\]

	\begin{block}{Numerical requirement}
		\[
			\boxed{n_{\text{Lemma A}} \approx 2.6\times 10^4}
		\]
	\end{block}

	\note{This is already quite large, but still the mildest of the three requirements.}
\end{frame}

% -----------------------
\begin{frame}{Lemma B: compressed operator concentration}
	\small
	The bound was
	\[
		\mathbb P\!\left(\|H-\mathbb EH\|_{\max}\ge \varepsilon\right)
		\le 2d^2\exp\!\left(-\frac{n\varepsilon^2}{32\kappa^2}\right).
	\]

	To make this at most \(\delta\), it is enough that
	\[
		2d^2\exp\!\left(-\frac{n\varepsilon^2}{32\kappa^2}\right)\le \delta,
	\]
	so
	\[
		\boxed{
			n \ge \frac{32\kappa^2}{\varepsilon^2}\log\!\left(\frac{2d^2}{\delta}\right).
		}
	\]

	Plugging in \(\kappa=\frac32\), \(\varepsilon=0.05\), \(d=9\), \(\delta=0.05\):
	\[
		n \ge \frac{32(3/2)^2}{0.05^2}\log\!\left(\frac{2\cdot 9^2}{0.05}\right)
		\approx 232{,}800.
	\]

	\begin{block}{Numerical requirement}
		\[
			\boxed{n_{\text{Lemma B}} \approx 2.3\times 10^5}
		\]
	\end{block}

	\note{This is the most expensive bound because McDiarmid uses only bounded differences and we are controlling all entries uniformly.}
\end{frame}

% -----------------------
\begin{frame}{Lemma C: approximate eigen-action}
	\small
	The bound was
	\[
		\mathbb P\!\left(\|A\Phi-\Phi\Lambda^{(n)}\|_{\max}\ge \varepsilon\right)
		\le 2nd\exp\!\left(-\frac{n\varepsilon^2}{8\kappa^2}\right).
	\]

	To make this at most \(\delta\), we need
	\[
		2nd\exp\!\left(-\frac{n\varepsilon^2}{8\kappa^2}\right)\le \delta.
	\]

	This is implicit because \(n\) appears inside the logarithm. Rearranging:
	\[
		\boxed{
			n \ge \frac{8\kappa^2}{\varepsilon^2}\log\!\left(\frac{2nd}{\delta}\right).
		}
	\]

	For \(\kappa=\frac32\), \(\varepsilon=0.05\), \(d=9\), \(\delta=0.05\), solving this fixed-point inequality gives
	\[
		n \approx 126{,}994.
	\]

	\begin{block}{Numerical requirement}
		\[
			\boxed{n_{\text{Lemma C}} \approx 1.3\times 10^5}
		\]
	\end{block}

	\note{The prefactor \(nd\) comes from a union bound over all rows and low-mode columns, which makes the estimate noticeably looser.}
\end{frame}

% -----------------------
\begin{frame}{Numerical values}
	\small
	\begin{table}
		\centering
		\begin{tabular}{l c}
			\hline
			\textbf{Lemma}                             & \textbf{Required \(n\) for \(\varepsilon=0.05,\ \delta=0.05,\ K=4\)} \\
			\hline
			Lemma A: Gram concentration                & \( \approx 25{,}867\)                                                \\
			Lemma B: compressed operator concentration & \( \approx 232{,}800\)                                               \\
			Lemma C: approximate eigen-action          & \( \approx 126{,}994\)                                               \\
			\hline
		\end{tabular}
	\end{table}

	\begin{block}{Takeaway}
		Even for the first \(K=4\) modes and a modest target \(\varepsilon=0.05\), the current worst-case concentration bounds ask for
		\[
			n \sim 10^4 \text{ to } 10^5.
		\]
	\end{block}

	\note{This is the slide to discuss with the advisor: rigorous, but numerically pessimistic.}
\end{frame}

% -----------------------
\begin{frame}{What this means}
	\small
	\begin{block}{Interpretation}
		\begin{itemize}
			\item The constants are \textbf{very conservative}.
			\item They are good for:
			      \begin{itemize}
				      \item consistency as \(n\to\infty\),
				      \item showing the right qualitative dependence on \(\varepsilon,\delta,K\),
				      \item justifying the continuum-to-finite-sample picture.
			      \end{itemize}
			\item They might \alert{not} be sharp enough to predict realistic sample sizes.
		\end{itemize}
	\end{block}

	\begin{block}{Improvements}
		To get meaningful finite-\(n\) rates, we probably need sharper tools:
		\begin{itemize}
			% \item control the compressed \(d\times d\) operator directly,
			\item use variance-aware bounds (Bernstein / matrix concentration),
			\item avoid raw union bounds over all rows.
		\end{itemize}
	\end{block}

	\note{The message is not “the theory fails”; it is “the first-pass concentration proof is too crude for practical \(n\).”}
\end{frame}

% -----------------------
\begin{frame}{If we relax the tolerance}
	\small
	Keeping \(K=4\), \(d=9\), \(\delta=0.05\), \(\kappa=\tfrac32\), but changing \(\varepsilon\):

	\begin{table}
		\centering
		\begin{tabular}{c c c c}
			\hline
			\(\varepsilon\) & Lemma A                    & Lemma C                    & Lemma B                     \\
			\hline
			0.10            & \(\approx 6.5\times 10^3\) & \(\approx 2.9\times 10^4\) & \(\approx 5.8\times 10^4\)  \\
			0.05            & \(\approx 2.6\times 10^4\) & \(\approx 1.3\times 10^5\) & \(\approx 2.3\times 10^5\)  \\
			0.02            & \(\approx 1.6\times 10^5\) & \(\approx 8.8\times 10^5\) & \(\approx 1.45\times 10^6\) \\
			\hline
		\end{tabular}
	\end{table}

	\begin{block}{Pattern}
		The sample size grows essentially like \(1/\varepsilon^2\), so pushing to small error rapidly becomes expensive.
	\end{block}

	\note{This makes the tradeoff visually obvious: the issue is not just the constants, but also the \(1/\varepsilon^2\) scaling.}
\end{frame}

\end{document}

\end{document}
